\notepage{71}{Second-order energy and the validity criterion}
The second-order energy correction is
\begin{align*}
 E_n^2&=\mel{n^0}{H'}{n^1}\\
 &=\sum_{m\ne n}
 \frac{\mel{n^0}{H'}{m^0}\mel{m^0}{H'}{n^0}}{E_n^0-E_m^0}\\
 &=\boxed{\displaystyle\sum_{m\ne n}
 \frac{\abs{\mel{n^0}{H'}{m^0}}^2}{E_n^0-E_m^0}},
\end{align*}
where Hermiticity of $H'$ is used in the last step.
\smallheading{When is the expansion useful?}
A necessary requirement is that the first-order state correction be
small compared with the unperturbed state. In particular, each mixing
coefficient should satisfy
\[
 \abs{\frac{\mel{m^0}{H'}{n^0}}{E_n^0-E_m^0}}\ll1\qquad(m\ne n).
\]
This depends on three things:
\begin{enumerate}
 \item The overall magnitude of $H'$.
 \item The matrix element $\mel{m^0}{H'}{n^0}$, which also depends on
       the nature of the two states.
 \item The energy spacing $\abs{E_n^0-E_m^0}$.
\end{enumerate}
If $E_n^0=E_m^0$, the nondegenerate expression has a zero denominator.
One must use degenerate perturbation theory instead.
\editorial{Small individual coefficients are a useful necessary check;
the norm of their combined correction must also remain small.}
\smallheading{Exercise retained from the notes}
Consider a charged harmonic oscillator in a uniform electric field
$\mathcal E$:
\[
 H=\frac{P^2}{2m}+\frac12m\omega^2X^2-q\mathcal E X.
\]
Solve the problem both by first-order perturbation theory and exactly.
The handwritten reading reminder refers to textbook pages 454--456;
the textbook and edition are not identified on this page.
